\section{Discussion}
\label{sec:discussion}

The main takeaway is deliberately modest: the experiments do not produce a
single winner between diffusion and autoregression. The ranking changes with
the sampling budget and with the parser constraints of the benchmark. Greedy
$n=1$ accuracy, estimated stochastic pass@1, low-$k$ sampling, and high-$k$
pass@k each reward a slightly different kind of behavior, especially when a
valid answer must satisfy several constraints at once.

Countdown-cd4 shows this most clearly. In the base condition, LLaDA-Base and
Dream-Base lead at estimated stochastic pass@1. By pass@128, Qwen-Base and
Llama-Base have overtaken them. The paired bootstrap comparisons support both
sides of the LLaDA--Qwen reversal. If only a single top-1 score were reported,
the larger gains from repeated Qwen and Llama sampling would be easy to miss.

Trip Planning follows the same broad theme, but with more noise. Dream is
strong in deterministic and low-$k$ settings, while flat-prompt Llama-Instruct
gives the strongest observed high-$k$ result. The caveat is important: many
candidate itineraries fail before semantic scoring. The benchmark therefore
mixes planning quality with the ability to produce the expected itinerary
format, which makes the high-$k$ differences less clean than on Countdown.

GSM8K plays a different role. The main runs approach saturation at high $k$, so
small ranking differences there are weak evidence about generation interfaces.
In this thesis, GSM8K is more useful as a check on prompt format and answer
extraction, especially the boxed-answer requirement in the primary strict
scorer.

The original motivation came from Yue et al.'s large-$k$ result: models that are
weaker at pass@1 can become stronger when many samples are inspected
\cite{yue2025rlvr}. Under the tested settings, LLaDA and Dream do not show the
analogous large-$k$ coverage effect. Qwen and Llama often gain more from
repeated stochastic sampling. That pattern is suggestive, but it should not be
over-read as a pure architectural conclusion, since prompt interface,
post-training, tokenizer behavior, and backend also vary.

Finally, the per-question and oracle-union analyses show overlap without
identity. Some questions are solved by one model family and not another. That is
useful evidence of complementarity, but only as an upper bound: the oracle
assumes correctness is known at selection time. A practical system would need a
verifier, router, selector, or parser-aware repair step to choose among
candidates without seeing the gold answer.
